\chapter{Yêu cầu về hệ thống}

\section{Yêu cầu chức năng}

\subsection{User App}
Các chức năng chính bao gồm:
\begin{enumerate}
    \item \textbf{Onboarding \& Cá nhân hóa Hồ sơ}
    \begin{itemize}
        \item \textbf{Mục đích:} Tạo trải nghiệm ban đầu liền mạch và xây dựng hồ sơ số để cá nhân hóa.
        \item \textbf{Chức năng cốt lõi}:
        \begin{itemize}
            \item \textbf{Đăng ký:} Tạo tài khoản qua \textbf{username/password} hoặc \textbf{Google/Facebook}.
            \item  \textbf{Khảo sát: } Thu thập nhanh nhu cầu ban đầu.
            \item \textbf{Khởi tạo hồ sơ: }Tự động tạo hồ sơ và đề xuất dịch vụ phù hợp.
        \end{itemize}
    \end{itemize}
    \item \textbf{Chatbot AI Tư vấn \& Gợi ý Dịch vụ}
    \begin{itemize}
        \item \textbf{Mục đích: }Trợ lý ảo 24/7 giúp người dùng xác định nhu cầu qua hội thoại.
        \item \textbf{Chức năng cốt lõi:}
        \begin{itemize}
            \item \textbf{Xử lý ngôn ngữ tự nhiên (NLP):} Hiểu và phân tích yêu cầu người dùng.
            \item \textbf{Tư vấn theo ngữ cảnh:} Đặt câu hỏi làm rõ nhu cầu.
            \item  \textbf{Đề xuất thông minh: }Gợi ý dịch vụ kèm thông tin chính và điều hướng đặt lịch.
        \end{itemize}
    \end{itemize}
    \item \textbf{Tìm kiếm \& Lọc Nâng cao}
    \begin{itemize}
        \item \textbf{Mục đích:} Cho phép tìm kiếm và lọc dịch vụ chính xác theo nhiều tiêu chí.
        \item \textbf{Chức năng cốt lõi: } 
        \begin{itemize}
            \item \textbf{Thanh tìm kiếm thông minh:} Tìm theo từ khóa, địa điểm, tên đối tác và tự động gợi ý.
            \item  \textbf{Bộ lọc đa chiều:} Lọc theo giá, đánh giá, vị trí và thuộc tính dịch vụ.
            \item \textbf{Hiển thị bản đồ:} Hiển thị vị trí nhà cung cấp theo khu vực đang xem.
        \end{itemize}
    \end{itemize}
    \item \textbf{Đặt lịch hẹn (Booking)}
    \begin{itemize}
        \item \textbf{Mục đích:} Số hóa quy trình đặt lịch để tăng chính xác và giảm sai sót.
        \item \textbf{Chức năng cốt lõi:}
        \begin{itemize}
            \item \textbf{Lịch hẹn thời gian thực:} Hiển thị khung giờ trống để người dùng chọn trực tiếp.
            \item \textbf{Tự động hóa thông báo: }Gửi xác nhận và nhắc lịch qua email/thông báo đẩy.
            \item \textbf{Quản lý tập trung: }Cho phép xem, đổi hoặc hủy lịch đã đặt.
        \end{itemize}
    \end{itemize}
    \item \textbf{Thanh toán \& Đánh giá}
    \begin{itemize}
        \item \textbf{Mục đích:} Đảm bảo giao dịch an toàn, minh bạch và xây dựng hệ thống phản hồi tin cậy.
        \item \textbf{Chức năng cốt lõi:} 
        \begin{itemize}
            \item \textbf{Tích hợp đa kênh thanh toán:} Hỗ trợ thẻ tín dụng, ví điện tử và cổng thanh toán uy tín.
            \item \textbf{Hệ thống đánh giá xác thực:} Chỉ người dùng đã hoàn thành dịch vụ mới được đánh giá.
            \item \textbf{Minh bạch thông tin:} Công khai đánh giá để hỗ trợ quyết định của người dùng.
        \end{itemize}
    \end{itemize}
    \item \textbf{AI trong Quản trị \& Tương tác thông minh}
    \begin{itemize}
        \item \textbf{Mục đích:} Tối ưu vận hành và tăng gắn kết người dùng bằng tương tác cá nhân hóa.
        \item \textbf{Chức năng cốt lõi}:
        \begin{itemize}
            \item \textbf{Dashboard phân tích:} Cung cấp báo cáo và biểu đồ về hiệu suất, hành vi người dùng.
            \item \textbf{Phân khúc khách hàng tự động:} AI nhóm người dùng theo hành vi cho chiến dịch phù hợp.
            \item \textbf{Giao tiếp chủ động:} Gửi thông báo, email marketing và ưu đãi cá nhân hóa đúng thời điểm.
        \end{itemize}
    \end{itemize}
\end{enumerate}

\subsection{Partner Portal}
Các chức năng chính bao gồm:
\begin{enumerate}
    \item \textbf{Đăng ký \& Hồ sơ}: Biểu mẫu đa bước, cổng tải tài liệu pháp lý, xác nhận email tự động.
    \item \textbf{Cổng Chờ duyệt}: Theo dõi trạng thái hồ sơ, kênh giao tiếp 2 chiều với Admin.
    \item \textbf{Quản lý Lịch hẹn}: Dashboard lịch hẹn, xác nhận/từ chối, calendar view ngày/tuần/tháng.
    \item \textbf{Quản lý Dịch vụ}: CRUD dịch vụ (tên, mô tả, giá, hình ảnh), gán nhân viên, cập nhật realtime.
    \item \textbf{Quản lý Nhân viên}: Tạo hồ sơ, phân quyền theo vai trò, tài khoản cá nhân hóa.
    \item \textbf{Khuyến mãi}: Tạo đa dạng (%, tiền, mua X tặng Y), điều kiện linh hoạt, hiển thị nổi bật.
    \item \textbf{Báo cáo}: Dashboard KPIs, thống kê dịch vụ/khung giờ, xuất Excel/CSV.
\end{enumerate}

\subsection{Admin Dashboard}
Các chức năng chính bao gồm:
\begin{enumerate}
    \item \textbf{Xác thực \& Quản lý Đối tác (Partner Verification \& Management)}
    \begin{itemize}
        \item \textbf{Mục đích:} Đảm bảo chất lượng và tính hợp pháp của toàn bộ đối tác trên nền tảng thông qua một quy trình xét duyệt và quản lý chuyên nghiệp.
        \item \textbf{Chức năng cốt lõi:}
        \begin{itemize}
            \item \textbf{Hệ thống xét duyệt tập trung:} Hiển thị tất cả hồ sơ đăng ký mới tại một nơi duy nhất để admin dễ dàng xem xét thông tin và các tài liệu pháp lý đính kèm.
            \item \textbf{Quy trình xử lý 3 chiều:} Cung cấp các hành động rõ ràng: Phê duyệt (mở khóa toàn bộ tính năng cho đối tác), Từ chối (kèm lý do), hoặc Yêu cầu bổ sung (gửi thông báo về các mục cần hoàn thiện).
            \item \textbf{Quản lý đối tác hiện hữu:} Cho phép admin thực hiện các tác vụ quản trị như tạm khóa tài khoản, cập nhật thông tin, và giám sát hoạt động.
        \end{itemize}
    \end{itemize}
    \item \textbf{Quản lý Đánh giá \& Báo cáo (Review \& Report Moderation)}
    \begin{itemize}
        \item \textbf{Mục đích:} Xây dựng và duy trì một môi trường cộng đồng văn minh, công bằng bằng cách kiểm duyệt các đánh giá và xử lý các báo cáo vi phạm.
        \item \textbf{Chức năng cốt lõi:}
        \begin{itemize}
            \item \textbf{Kiểm duyệt nội dung:} Admin có thể duyệt hoặc loại bỏ các đánh giá chứa nội dung không phù hợp, spam, hoặc vi phạm chính sách.
            \item \textbf{Xử lý báo cáo:} Tiếp nhận và giải quyết các báo cáo vi phạm từ người dùng về đối tác (hoặc ngược lại), đảm bảo quyền lợi cho các bên.
        \end{itemize}
    \end{itemize} 
    \item \textbf{Quản lý Tài chính Toàn hệ thống (System-wide Financial Management)}
    \begin{itemize}
        \item \textbf{Mục đích:} Cung cấp một công cụ tổng thể để giám sát dòng tiền, quản lý hoa hồng, và tự động hóa quy trình đối soát, thanh toán cho đối tác.
        \item \textbf{Chức năng cốt lõi:}
        \begin{itemize}
            \item \textbf{Dashboard tài chính tổng quan:} Cung cấp cái nhìn toàn cảnh về các chỉ số tài chính quan trọng như tổng giá trị giao dịch (GMV) và doanh thu nền tảng.
            \item \textbf{Hệ thống đối soát tự động:} Tự động tổng hợp doanh thu và tính toán số tiền cần thanh toán cho mỗi đối tác theo chu kỳ đã định sẵn (ví dụ: hàng tháng).
            \item \textbf{Quản lý thanh toán:} Hỗ trợ thực hiện thanh toán hàng loạt và theo dõi trạng thái công nợ của từng đối tác một cách minh bạch.
        \end{itemize}
    \end{itemize} 
    % \item \textbf{Quản lý Mã giảm giá \& Chiến dịch Marketing (Voucher \& Campaign Management)}
    % \begin{itemize}
    %     \item \textbf{Mục đích:} Cung cấp một công cụ marketing tập trung để admin tạo và điều phối các chiến dịch quảng bá trên toàn hệ thống, thúc đẩy tăng trưởng chung.
    %     \item \textbf{Chức năng cốt lõi:}
    %     \begin{itemize}
    %         \item \textbf{Tạo chiến dịch tổng thể:} Cho phép khởi tạo các chiến dịch lớn theo chủ đề (ví dụ: "Chào hè", "Black Friday") và quản lý tập trung.
    %         \item \textbf{Phát hành mã giảm giá toàn hệ thống:} Tạo các mã voucher áp dụng chung cho nhiều đối tác, với các quy tắc linh hoạt (giới hạn lượt dùng, giá trị đơn tối thiểu, áp dụng cho người dùng mới...).
    %         \item \textbf{Theo dõi hiệu quả: }Đo lường hiệu suất của từng chiến dịch thông qua các chỉ số như số lượt sử dụng mã, doanh thu tạo ra.
    %     \end{itemize}
    % \end{itemize}
    % \item \textbf{Trung tâm Hỗ trợ \& Ticket (Support Center \& Ticketing System)}
    % \begin{itemize}
    %     \item \textbf{Mục đích:} Chuẩn hóa và tối ưu hóa quy trình hỗ trợ khách hàng, đảm bảo mọi yêu cầu, khiếu nại đều được ghi nhận, phân công và xử lý triệt để.
    %     \item \textbf{Chức năng cốt lõi:}
    %     \begin{itemize}
    %         \item \textbf{Hệ thống gom ticket:} Tập hợp tất cả yêu cầu hỗ trợ từ người dùng và đối tác vào một hàng đợi duy nhất.
    %         \item \textbf{Theo dõi trạng thái: }Gán trạng thái cho mỗi ticket (Mới, Đang xử lý, Đã giải quyết) để dễ dàng theo dõi tiến độ.
    %         \item \textbf{Hỗ trợ trực tiếp qua Trò chuyện (Live Chat for Users): }Người dùng có thể bắt đầu một cuộc trò chuyện trực tiếp với đội ngũ hỗ trợ ngay trên ứng dụng để giải quyết các vấn đề cấp bách hoặc các câu hỏi nhanh.   
    %     \end{itemize}
    % \end{itemize}
\end{enumerate}


\section{Yêu cầu phi chức năng}
\begin{enumerate}
    \item \textbf{Hiệu suất (Performance):}
    \begin{itemize}
        \item \textbf{Thời gian phản hồi (Response Time):} Mọi tương tác trên ứng dụng (tải trang, tìm kiếm, lọc) phải hoàn thành dưới \textbf{2 giây}. API response time trung bình phải dưới \textbf{200ms} và tối đa (P95) không vượt \textbf{500ms}.
        \item \textbf{Tải trọng (Load):} Hệ thống phải có khả năng phục vụ đồng thời ít nhất \textbf{1.000 người dùng} (concurrent users) trong giờ cao điểm mà không suy giảm hiệu suất đáng kể (< 10\% tăng latency).
        \item \textbf{Throughput:} Hệ thống backend phải xử lý được tối thiểu \textbf{500 requests/giây} cho các API chính (tìm kiếm, đặt lịch, thanh toán).
        \item \textbf{Hiệu suất AI Chatbot:} Thời gian phản hồi của chatbot AI (bao gồm truy xuất RAG và sinh văn bản) không vượt quá \textbf{15 giây} cho mỗi lượt hội thoại. Tỷ lệ phản hồi chính xác (relevance) đạt tối thiểu \textbf{85\%}.
        \item \textbf{Tối ưu hóa mobile:} Ứng dụng di động phải khởi động (cold start) trong vòng \textbf{3 giây} trên các thiết bị tầm trung. Kích thước bộ nhớ (RAM) sử dụng không vượt \textbf{150MB} trong quá trình hoạt động bình thường.
    \end{itemize}
    \item \textbf{Bảo mật (Security):}
    \begin{itemize}
        \item \textbf{Mã hóa dữ liệu:} Toàn bộ dữ liệu phải được mã hóa khi truyền (\textbf{TLS 1.2+}) và khi lưu trữ (\textbf{AES-256} encryption at rest), đặc biệt là thông tin cá nhân và lịch sử sức khỏe của người dùng.
        \item \textbf{Xác thực \& Phân quyền:} Sử dụng JWT (JSON Web Token) với thời gian hết hạn ngắn (access token \textbf{15 phút}, refresh token \textbf{7 ngày}). Áp dụng chặt chẽ cơ chế phân quyền theo vai trò (\textbf{RBAC}) với tối thiểu 4 vai trò: User, Partner, Staff, Admin.
        \item \textbf{Bảo vệ API:} Triển khai \textbf{rate limiting} (tối đa 100 requests/phút/user), \textbf{input validation} trên mọi endpoint, và bảo vệ chống các cuộc tấn công phổ biến (SQL Injection, XSS, CSRF).
        \item \textbf{Tuân thủ:} Hệ thống phải tuân thủ các tiêu chuẩn bảo mật \textbf{OWASP Top 10} và các quy định về bảo vệ dữ liệu cá nhân tại Việt Nam (Nghị định 13/2023/NĐ-CP).
    \end{itemize}
    \item \textbf{Độ sẵn sàng (Availability):}
    \begin{itemize}
        \item \textbf{Uptime:} Hệ thống phải đảm bảo thời gian hoạt động \textbf{99.9\%} (tức là thời gian downtime không quá \textbf{43 phút/tháng}), không bao gồm thời gian bảo trì theo lịch.
        \item \textbf{Sao lưu dữ liệu:} Cơ sở dữ liệu phải được sao lưu tự động (\textbf{automated backup}) ít nhất \textbf{1 lần/ngày}, với khả năng phục hồi (RTO) trong vòng \textbf{1 giờ} và điểm phục hồi (RPO) không quá \textbf{4 giờ}.
        \item \textbf{Graceful Degradation:} Khi một thành phần phụ (ví dụ: AI Service, Notification Service) gặp sự cố, các chức năng cốt lõi (đặt lịch, thanh toán) vẫn phải hoạt động bình thường.
    \end{itemize}
    \item \textbf{Khả năng mở rộng (Scalability):}
    \begin{itemize}
        \item \textbf{Horizontal Scaling:} Kiến trúc Microservices cho phép mở rộng theo chiều ngang, mỗi service có thể scale độc lập dựa trên nhu cầu thực tế (ví dụ: AI Service có thể scale riêng khi lượng truy vấn chatbot tăng).
        \item \textbf{Database Scaling:} PostgreSQL phải hỗ trợ \textbf{read replicas} để phân tải truy vấn đọc và \textbf{connection pooling} để quản lý kết nối hiệu quả.
        \item \textbf{Mục tiêu dài hạn:} Hệ thống phải sẵn sàng phục vụ tối thiểu \textbf{100.000 người dùng} và \textbf{1.000 đối tác} trong vòng 12 tháng sau khi ra mắt mà không cần tái cấu trúc kiến trúc.
    \end{itemize}
    \item \textbf{Tính khả dụng \& Trải nghiệm người dùng (Usability \& UX):}
    \begin{itemize}
        \item \textbf{Tương thích thiết bị:} Ứng dụng phải hoạt động mượt mà trên các thiết bị iOS (từ iOS 14+) và Android (từ Android 8+) phổ biến trên thị trường.
        \item \textbf{Nhất quán giao diện:} Giao diện và trải nghiệm người dùng phải nhất quán giữa các màn hình và nền tảng (User App, Partner Portal, Admin Dashboard), tuân thủ Material Design Guidelines.
        \item \textbf{Hỗ trợ người dùng:} Cung cấp tài liệu hướng dẫn (FAQ), tooltip trên các chức năng phức tạp, và các kênh hỗ trợ dễ tiếp cận ngay trong ứng dụng/portal.
        %\item \textbf{Accessibility:} Giao diện phải đảm bảo tỷ lệ tương phản tối thiểu \textbf{4.5:1} cho văn bản, hỗ trợ cỡ chữ linh hoạt (text scaling), và tương thích với các công cụ hỗ trợ người khuyết tật (screen reader).
        %\item \textbf{Đa ngôn ngữ:} Hệ thống phải hỗ trợ tối thiểu \textbf{tiếng Việt} (mặc định) và \textbf{tiếng Anh}, với khả năng mở rộng thêm ngôn ngữ trong tương lai.
    \end{itemize}
    \item \textbf{Khả năng bảo trì (Maintainability):}
    \begin{itemize}
        \item \textbf{Code Quality:} Mã nguồn phải được viết sạch, tuân thủ coding convention (ESLint/Prettier cho TypeScript, Dart Analyzer cho Flutter), có tài liệu và comment rõ ràng.
        \item \textbf{Kiến trúc module hóa:} Áp dụng Clean Architecture với phân tách rõ ràng giữa các layer (Presentation, Domain, Data), cho phép thay đổi một layer mà không ảnh hưởng đến các layer khác.
        %\item \textbf{CI/CD:} Triển khai pipeline tự động hóa cho việc build, test, và deploy. Mỗi commit phải qua ít nhất \textbf{unit test} và \textbf{lint check} trước khi merge.
        \item \textbf{Logging \& Monitoring:} Hệ thống phải tích hợp logging có cấu trúc (structured logging) và monitoring dashboard để theo dõi health check, error rate, và các metrics quan trọng theo thời gian thực.
    \end{itemize}

    \item \textbf{Độ tin cậy (Reliability):}
    \begin{itemize}
        \item \textbf{Tính toàn vẹn dữ liệu:} Mọi giao dịch quan trọng (đặt lịch, thanh toán) phải đảm bảo tính ACID. Sử dụng cơ chế \textbf{idempotency} cho các API thanh toán để tránh trùng lặp giao dịch.
        \item \textbf{Error Handling:} Hệ thống phải có cơ chế xử lý lỗi thống nhất, trả về mã lỗi rõ ràng (theo chuẩn HTTP status codes) và thông báo thân thiện với người dùng. Tỷ lệ lỗi hệ thống (5xx errors) không vượt quá \textbf{0.1\%} tổng số requests.
        \item \textbf{Retry \& Circuit Breaker:} Các lời gọi đến dịch vụ bên thứ ba (Payment Gateway, Push Notification) phải có cơ chế retry tự động (tối đa 3 lần) và circuit breaker để ngăn cascading failure.
    \end{itemize}
\end{enumerate}

